% !TEX root = ../main.tex
\chapter{Tensors and tensor fields}\label{ch:c05}
\begin{paracol}{2}

\begin{sources}
\begin{tabularx}{\linewidth}{@{}l l L@{}}
\toprule
\textbf{Video} & \textbf{Length} & \textbf{Content}\\
\midrule
\seg{5}{1}{Lecture 5: Tensors and tensor fields (Tensoren und Tensorfelder)} & 1:53:51 & Recap of tangent vectors in a chart and of their
transformation law; defining a vector by its components in a chart; constructing components from a given object
(first derivative of a curve works, second derivative does not); vector fields; $(r,s)$ tensor fields and their transformation
law; the Kronecker delta and the inverse Jacobian matrices; examples: the gradient is a covector field, the matrix of second
derivatives is not a tensor field.\\
\bottomrule
\end{tabularx}

\smallskip
\textbf{Transcript} \file{materials/transcripts/C05-V01.txt} (German YouTube auto-captions).
\end{sources}

\begin{reading}[Where this lecture sits in the other notes]
\emph{Geometric Anatomy}: lectures 8 and 11 (tensor space theory I and II: tensors as multilinear maps, components,
change of basis), lecture 9 (differential structures, vector fields), and lecture 25 (covariant derivatives, which
repair the non-tensorial second derivative announced here).
\emph{Gravity and Light} (TP-1 version): lectures 5 and 6 (tensors and tensor fields).
\end{reading}

% ==================================================================
\section{Why tensors?}\label{sec:tf-motivation}

``What is a tensor?'' The answer is the content of the lecture. What matters for you: physics teems with tensors, you
cannot avoid them, and you have already met many without knowing it; most of what is called a tensor is not one, only in
special cases \ts{5}{1}{0:04}. A tensor is a generalisation of the concept of a vector: ``in quotation marks, one can argue
about this characterisation, it is only to whet the appetite: \emph{vectors of higher rank}'' \ts{5}{1}{1:10}.
The name comes from \emph{tension} (Spannung): the stress tensor of a medium was one of the first physical applications
(the principal-axes trick for a rank-two tensor) \ts{5}{1}{1:08:15}.

% ==================================================================
\section{Recap: vectors on the chart level}\label{sec:tf-recap}

Vectors are defined already at the manifold level, as directional derivatives \ts{5}{1}{2:11}. Recall:
$(M,\mathcal O,\mathcal A)$ is a smooth manifold; the tangent vector $X_{\gamma,p}\colon C^\infty(M)\to\R$,
$f\mapsto(f\circ\gamma)'(\lambda_0)$, is a \emph{map}, not a collection of numbers (a directional derivative operator, since it
acts on $f$ before producing anything) \ts{5}{1}{3:34}.
After the choice of a chart $(U,x)$ with $p\in U$ it can be written
\[
  X_{\gamma,p}=\dot\gamma^{\,j}_{(x)}(\lambda_0)\left(\frac{\partial}{\partial x^j}\right)_{\!p},
  \qquad \dot\gamma^{\,j}_{(x)}=(x^j\circ\gamma)',
\]
a sum over $j=1,\dots,d$ of the components times the chart-induced basis of the vector space $T_pM$, at the point $p$ only
(the basis at another point lives in another tangent space) \ts{5}{1}{6:50}.
Once more the warning \ts{5}{1}{11:43}: $(\partial/\partial x^j)_p$ is \emph{not} a partial derivative; acting on $f$ it is
defined as $\partial_j(f\circ x^{-1})(x(p))$, the derivative of $f\circ x^{-1}\colon\R^d\to\R$ in its $j$-th slot; whenever you
meet the symbol, first convert it to this and then compute \ts{5}{1}{11:43}.

In a different chart $(V,y)$ the same abstract vector has other components and the basis changes
(\cref{ch:c04}):
\begin{equation}\label{eq:tf-vector-trafo}
  \dot\gamma^{\,i}_{(y)}\;=\;\frac{\partial y^i}{\partial x^m}\Big|_{x(p)}\;\dot\gamma^{\,m}_{(x)},
\end{equation}
where $\partial y^i/\partial x^m|_{x(p)}:=\partial_m(y^i\circ x^{-1})(x(p))$ is the derivative of the $i$-th component of the
transition map $y\circ x^{-1}$ \ts{5}{1}{14:30}.
This is the \textbf{transformation law for vector components} (Transformationsformel für Vektorkomponenten), with the
summation convention: an index repeated once up and once down is summed. ``There is nothing simpler: if you have the
components $2,3,7$ in one chart, you get those in the new chart through this formula''
\ts{5}{1}{17:59}.
\begin{coursenote}[What $y^i$ means in $\partial y^i/\partial x^m$]
The symbol $\partial y^i/\partial x^m|_{x(p)}$ is ambiguous at first sight, since $y^i$ is a function on the chart domain
$V\subseteq M$ and not on $\R^d$: it is the $i$-th chart-coordinate function. Smoothness: $y^i\circ y^{-1}$ is the
projection of $\R^d$ to the $i$-th entry, hence smooth, so $y^i\in C^\infty$ \emph{provided} $y$ belongs to the smooth atlas. The
object is defined as above: one needs the explicit transition map, then differentiates its $i$-th component in the $m$-th
entry \ts{5}{1}{18:41}.
\end{coursenote}
The calculations of the last lecture should be studied ``until the last equals sign is understood'': the technical key is the
chain rule. ``This is not the children's version: there will probably be no other definition until the end of your
studies and your Nobel Prize'' \ts{5}{1}{22:32}.

% ==================================================================
\section{Vector definition on the chart level}\label{sec:tf-vector-def}

Formula \eqref{eq:tf-vector-trafo} can be used as a \emph{definition} \ts{5}{1}{23:13}. We want a
vector at the point $p$, without a curve. Choose a chart $(U,x)$ containing $p$, and \emph{prescribe} components, ``just as
we like'': e.g.\ $X^1(p)=\ldots$, $X^2(p)=7$, $X^3(p)=-5$ (the numbers are in the transcript as ``$X^1,X^2=7,X^3=-5$''); the
object $X^i(p)\,(\partial/\partial x^i)_p$ is a vector in $T_pM$ for any choice of numbers
\ts{5}{1}{24:16}.

\begin{definition}[New concepts: vector defined by components, chart-independence]\label{def:tf-vector}
A vector at $p$ is given by its components $X^i_{(x)}(p)$ with respect to a chart $(U,x)$, \emph{together with the demand}
that the components with respect to any other chart $(V,y)$ with $p\in V$ are
\[
  X^i_{(y)}(p)\;=\;\frac{\partial y^i}{\partial x^m}\Big|_{x(p)}\;X^m_{(x)}(p)
\]
\ts{5}{1}{26:21}
(the little label $(x)$ or $(y)$ under the symbol records the chart to which the components refer; later one drops it or uses
a tilde when changing chart).
\end{definition}
So the logic is now the one of school: one gives a vector by components, but to make it a real tangent vector one must
\emph{require} this transformation behaviour. When one wishes to change chart one uses the transition map. Before, we
could \emph{derive} the formula from the abstract definition on the manifold; now we start from the components. What was said in
school and experimental physics (``the vector with components $3,7,5$'') is not wrong: a silent bracket opens ``with
respect to a chosen coordinate system'', and if one changes it one uses this formula \ts{5}{1}{29:11}.
Everything one knew about vectors is only made precise \ts{5}{1}{30:17}.

\subsection{Constructing components from a given object}\label{ssec:tf-construct}

Usually one has no occasion to write down such a definition; more often one has already a mathematical object in a chart
(a curve, a function), and constructs ``components'' from it. One cannot simply call them vector components: knowing how the
object they come from transforms, one can \emph{derive} how the constructed components transform and then check whether this
agrees with the vector law. It is as common that it does not as that it does \ts{5}{1}{30:17}.
Care is needed when one ``defines'' components by constructing them from a mathematical object (curves, functions, \dots) in a chart;
the transformation law of the constructed components may deviate from the vector law \ts{5}{1}{32:02}.

\begin{example}[First derivative of a curve: components of a vector]\label{ex:tf-first}
Let $q^i_{(x)}=x^i\circ\gamma$ be the chart representation of a curve and define $\dot q^i_{(x)}(\lambda):=(q^i_{(x)})'(\lambda)$.
In another chart, $\dot q^i_{(y)}=(y^i\circ\gamma)'$, and by the chain rule applied to $q_{(y)}=(y\circ x^{-1})\circ q_{(x)}$,
\[
  \dot q^i_{(y)}(\lambda)=\frac{\partial y^i}{\partial x^m}\Big|_{x(\gamma(\lambda))}\dot q^m_{(x)}(\lambda),
\]
the vector law: they are genuine components of a vector (as shown earlier; this is problem sheet~4, exercise~2, to be computed
by hand) \ts{5}{1}{34:19}.
In the chart $x$ one may draw a true arrow with these components, ``on the manifold it is a directional derivative''
\ts{5}{1}{38:11}.
\end{example}

\begin{example}[Counterexample: the second derivative of a curve]\label{ex:tf-second}
Now take the second derivative, $\ddot q^i_{(x)}(\lambda):=(q^i_{(x)})''(\lambda)$: these are the ``components of the acceleration
vector''. Is that a vector? \ts{5}{1}{38:53}.
Differentiate $q^i_{(y)}=(y^i\circ x^{-1})\circ q_{(x)}$ twice (careful: do not substitute too early; one needs both product
and chain rule):
\[
  \ddot q^i_{(y)}
  =\frac{\partial y^i}{\partial x^m}\Big|_{x(\gamma(\lambda))}\ddot q^m_{(x)}
  \;+\;\frac{\partial^2 y^i}{\partial x^m\partial x^n}\Big|_{x(\gamma(\lambda))}\dot q^m_{(x)}\dot q^n_{(x)} .
\]
\ts{5}{1}{41:02}.
Were it not for the second term, the $\ddot q^i$ would transform as vector components. But the extra term, \emph{quadratic}
in the velocity components, does not vanish in general and destroys the tensor property: the second derivative of the chart
representation of a curve is \textbf{not} a vector. Neither are ``position vectors'' (arbitrary chart maps, not only linear
ones), nor the second derivative of the track \ts{5}{1}{43:28}.
\end{example}
\begin{aside}[Why Newton's law cannot be $m\ddot x=F$]
Hence a theorist cannot write $m\ddot{\vec x}=\vec F$ as Newton's second law: with a vector on the right-hand side the left
cannot be a vector. Problem sheet~4, exercise~2 checks it by hand. A student had asked by e-mail whom to believe, given
that the forbidden word ``position vector'' occurs in a book: ``believe nobody, compute it''
\ts{5}{1}{44:29}.
\end{aside}
Next lecture a \emph{new structure} on the manifold, a \textbf{linear connection} (also indexed quantities, transforming
differently), will supply a term whose transformation behaviour exactly cancels the extra term, so that together one gets a
genuine acceleration vector; Newton's law speaks of acceleration and force, so we must construct the acceleration correctly:
``the second derivative of the chart image of the curve is certainly not it; we will repair it'' \ts{5}{1}{46:12}.

\paragraph{Notation for higher derivatives.} A student asks what $\partial^2/\partial x^i\partial x^j$ means. It is
$\partial_i\partial_j(f\circ x^{-1})(x(p))$: first the $j$-th then the $i$-th slot derivative, evaluated at $x(p)$; higher ones
analogously; evaluation always at the end, after all differentiations \ts{5}{1}{47:55}.

% ==================================================================
\section{Vector fields}\label{sec:tf-vector-fields}

After a break \ts{5}{1}{50:01}. So far vectors at a \emph{single} point $p$ (drawn as an arrow, which in a chart one may do)
\ts{5}{1}{50:22}.

\begin{definition}[New concepts: vector field]\label{def:tf-vectorfield}
A \textbf{vector field} (Vektorfeld) $X$ assigns to every point $p$ of $M$ (or of a subset) a vector $X_p\in T_pM$ in the tangent
space \emph{at that point}: ``one never assigns to a point a vector in a different tangent space''
\ts{5}{1}{51:04}. Dropping the index $p$ that recorded the base point,
\[
  X\colon p\mapsto X_p=X^i_{(x)}(p)\left(\frac{\partial}{\partial x^i}\right)_{\!p},
\]
and in a chart the components are \textbf{component functions} $X^i_{(x)}\colon x(U)\to\R$ (functions of the point, hence of
the chart coordinates) \ts{5}{1}{52:47}.
It is a vector field if and only if the components transform with the vector law
\begin{equation}\label{eq:tf-vf-trafo}
  X^i_{(y)}\bigl(y(p)\bigr)\;=\;\frac{\partial y^i}{\partial x^m}\Big|_{x(p)}\;X^m_{(x)}\bigl(x(p)\bigr),
\end{equation}
now for all points $p$ of the overlap \ts{5}{1}{58:21}.
It is called \textbf{smooth} if its component functions are smooth for every chart of the atlas (a standard analysis notion)
\ts{5}{1}{1:02:11}.
\end{definition}
\begin{intuition}[What the formalism now reflects]
In school physics one speaks of the velocity ``$2,3,4$'' and of ``$5,6,7$'' at another place; the formalism does not say that
they are attached to different points. Here the base point is recorded \ts{5}{1}{54:53}.
Smooth: the vector at each point varies in the same way and makes no jump, which is the intuition one always had
\ts{5}{1}{1:02:53}.
\end{intuition}
The transformation law \eqref{eq:tf-vf-trafo} has a new feature: the point of evaluation moves; it is $y(p)$ in the $y$-chart and
the image $x(p)$ of the same point in the $x$-chart \ts{5}{1}{59:45}. Specialised to one point, it is the law
for vectors.

\begin{example}[Tangent vector field of a curve]\label{ex:tf-curve-field}
If a curve is given, one has at each point of the curve a tangent vector; in the chart one may draw real arrows. It is a vector
field only on the \emph{subspace} (the image of the curve), not on the whole manifold; ``we can allow this: a vector field on a
sub-manifold''. Later, in spacetime, we shall also consider vector fields only on a hypersurface (a spatial slice)
\ts{5}{1}{55:57}.
\end{example}
Again there are the two standpoints: the \emph{definition} standpoint (prescribe components in one chart and demand the law
\eqref{eq:tf-vf-trafo}) and the \emph{construction} standpoint (generate components from a known object and derive whether they
obey the law) \ts{5}{1}{1:00:27}.

\begin{secondary}[The tangent bundle]
Where does $X$ ``land''? Not in a single tangent space, but in the disjoint union $TM$ of all tangent spaces, the \textbf{tangent
bundle} (Tangentialbündel); a vector field is then a \emph{section} of the bundle. This needs a bundle projection and is not
elaborated here: ``the intuition is right and what happens on the chart level is mathematically precise''
\ts{5}{1}{1:03:13}.
\end{secondary}
A student asks why vectors, which we first generated from curves, can now be defined without a curve. Answer: the tangent
space is defined by curves, but one can always find a curve whose tangent vector is a given vector (in fact many, with the
same result); so we are not dependent on curves any more, and we can regard vectors as independent objects, and if asked
for a concrete one, say ``take this curve''. ``In many books: a vector is what transforms like a vector; this is not very
helpful, as it already contains the transformation law; but now you know that components must be shown, either by theory or by
construction, to transform that way'' \ts{5}{1}{1:05:22}.

% ==================================================================
\section{Tensors and tensor fields}\label{sec:tf-tensors}

Now the actual topic \ts{5}{1}{1:07:26}. (``Tensor'': the word comes from \emph{tension}; see above.)
\ts{5}{1}{1:08:15}

\begin{definition}[New concepts: $(r,s)$ tensor field]\label{def:tf-tensorfield}
Let $r,s\in\mathbb N_0$. An \textbf{$(r,s)$ tensor field} $T$ (Tensorfeld) on $M$ (or on a subset) assigns to every point $p$ an
$(r,s)$ tensor $T_p$ at $p$ (nothing to do with the tangent space symbol $T_pM$; Schuller suggests one might better write
``Tennis'' instead of $T$!). The pair $(r,s)$ is the \textbf{valence} or rank; ``second-rank tensor'' alone is not precise,
since $2=2+0=1+1=0+2$ \ts{5}{1}{1:09:07}.
In a chart $(U,x)$ it has the form
\[
  T=T^{i_1\dots i_r}_{\;j_1\dots j_s}(x)\;
  \left(\frac{\partial}{\partial x^{i_1}}\right)\otimes\dots\otimes\left(\frac{\partial}{\partial x^{i_r}}\right)\otimes
  d x^{j_1}\otimes\dots\otimes d x^{j_s},
\]
with sums over all indices: $r$ indices up, $s$ indices down, each running from $1$ to $d$; so a tensor of valence $(r,s)$ has
$d^{\,r+s}$ components. The symbol $\otimes$ is the \textbf{tensor product}: a new operation (Schuller teases: ``I do not tell you
what it is; only the first factor will be explained, the second is written to marvel at'') which makes
$\{\partial_{i_1}\otimes\dots\otimes dx^{j_s}\}$ a basis of the space of $(r,s)$ tensors at $p$
\ts{5}{1}{1:10:51}.
\end{definition}
(The $d^{r+s}$ is stated in class as ``$d$ to the power $r+s$'' \ts{5}{1}{1:17:51}.) The lower indices cannot be fed to a vector
basis $\partial/\partial x$, which has an upper-type index; the dual basis $dx^j$ of the cotangent space (\cref{ch:c04}) takes
them: so vectors are the $(1,0)$ case with basis $\partial/\partial x^i$, covectors the $(0,1)$ case with basis $dx^j$
\ts{5}{1}{1:14:38}.

\begin{theorem}[Transformation law of tensor components]\label{thm:tf-tensor-trafo}
For the abstract object to be independent of the chart, the component functions in another chart $(V,y)$ must be
\[
  T^{i_1\dots i_r}_{(y)\;j_1\dots j_s}\bigl(y(p)\bigr)
  =\frac{\partial y^{i_1}}{\partial x^{m_1}}\Big|_{x(p)}\!\!\cdots\frac{\partial y^{i_r}}{\partial x^{m_r}}\Big|_{x(p)}\;
   \frac{\partial x^{n_1}}{\partial y^{j_1}}\Big|_{y(p)}\!\!\cdots\frac{\partial x^{n_s}}{\partial y^{j_s}}\Big|_{y(p)}\;
   T^{m_1\dots m_r}_{(x)\;n_1\dots n_s}\bigl(x(p)\bigr)
\]
(sum over $m_k,n_l$): each upper index acquires a factor $\partial y/\partial x$ exactly as a vector index, and each lower index
the \emph{inverse} factor $\partial x/\partial y$ ``in the literal sense: where the first has $y$ up and $x$ down, the second has
$x$ up and $y$ down'' \ts{5}{1}{1:16:47}.
\end{theorem}
``This formula one need not memorise; with it one can at most frighten small children. One merely adds further indices and does
the same again as for vectors; upper indices go with $\partial y/\partial x$, lower ones with the opposite factor''
\ts{5}{1}{1:21:18}.

\subsection{The inverse Jacobian and the Kronecker delta}\label{ssec:tf-delta}

The factor for lower indices is indeed the inverse of the one for upper indices \ts{5}{1}{1:22:12}.

\begin{definition}[New concepts: Kronecker delta]\label{def:tf-kronecker}
$\delta^a{}_b:=1$ if $a=b$, $0$ if $a\ne b$ (one index up, one down: it is a $(1,1)$ tensor,
not $\delta_{ab}$ or $\delta^{ab}$) \ts{5}{1}{1:26:05}.
\end{definition}
\begin{theorem}[The two Jacobians are inverse]\label{thm:tf-inverse}
$\dfrac{\partial y^a}{\partial x^m}\Big|_{x(p)}\dfrac{\partial x^m}{\partial y^b}\Big|_{y(p)}=\delta^a{}_b$: the matrices of
the two Jacobians are mutually inverse in the sense of linear algebra \ts{5}{1}{1:22:34}.
\end{theorem}
\begin{proof}
The index $m$ appears once up, once down, so it is summed. By the definitions the first factor is
$\partial_m(y^a\circ x^{-1})(x(p))$, the second is $\partial_b(x^m\circ y^{-1})(y(p))$. The multidimensional chain rule
\emph{backwards} (the sum over the repeated index is a chain rule for the composite $(y^a\circ x^{-1})\circ(x\circ y^{-1})$) gives
the derivative in the $b$-th direction of $y^a\circ x^{-1}\circ x\circ y^{-1}=y^a\circ y^{-1}$, the $a$-th component of the
identity of $\R^d$, which is $\delta^a{}_b$ \ts{5}{1}{1:22:55}.
\end{proof}
So, at a point, the product of the two Jacobian matrices is the unit matrix: the upper-index and lower-index transformations are
mutually inverse \ts{5}{1}{1:27:28}.

% ==================================================================
\section{Examples}\label{sec:tf-examples}

\begin{example}[Vector fields are $(1,0)$ tensor fields]\label{ex:tf-vf-tensor}
A vector field is a $(1,0)$ tensor field (the case $r=1,s=0$ of \cref{def:tf-tensorfield}) \ts{5}{1}{1:28:10}.
\end{example}

\begin{example}[The gradient is a $(0,1)$ tensor field]\label{ex:tf-gradient}
Let $f\colon M\to\R$ be smooth and $(U,x)$ a chart. Construct ``would-be'' components with a lower index,
\[
  (\mathrm{grad}f)_{(x)\,j}:=\partial_j\bigl(f\circ x^{-1}\bigr)
\]
(an index in the denominator counts as a lower index: ``the notation is chosen so that it fits'')
\ts{5}{1}{1:29:57}.
If these are the components of a tensor field, it can only be a $(0,1)$ field, having one lower index
\ts{5}{1}{1:32:29}. Check the transformation: in the chart $y$, inserting $x^{-1}\circ x$ and
using the chain rule,
\begin{align*}
  \partial_j\bigl(f\circ y^{-1}\bigr)\bigl(y(p)\bigr)
  &=\partial_j\bigl((f\circ x^{-1})\circ(x\circ y^{-1})\bigr)\bigl(y(p)\bigr)\\
  &=\partial_m\bigl(f\circ x^{-1}\bigr)\bigl(x(p)\bigr)\;\partial_j\bigl(x^m\circ y^{-1}\bigr)\bigl(y(p)\bigr),
\end{align*}
i.e.\ $(\mathrm{grad} f)_{(y)\,j}=\dfrac{\partial x^m}{\partial y^j}\,(\mathrm{grad}f)_{(x)\,m}$: exactly the law for a lower
index (the $y$ below, the $x$ above) \ts{5}{1}{1:33:10}.
So it is a $(0,1)$ tensor field (the first one shown), the \textbf{gradient} of $f$ (the covector field $df$; in the notation of
\cref{ch:c04}, $(df)_p=d_pf$). Note: the \emph{gradient is not a vector field}; it is a covector field
\ts{5}{1}{1:37:28}.
\end{example}

\begin{aside}[``The gradient is perpendicular to the level lines'']
The usual characterisation: draw the level lines $f=1$, $f=2$ in two dimensions, the gradient is the vector perpendicular to
them. ``Alarm bell'': to say ``perpendicular'' one needs a scalar product, and we have none on our tangent spaces. But the
gradient is a covector and \emph{can eat a vector}: take a tangent vector $X$ to a level line (the tangent of a curve of
constant $f$); then $\partial_jf\;X^j=X(f)=0$ with no scalar product needed (``one can show it is always zero''). The
usual claim ``the gradient is a vector perpendicular to the level lines'' needs a scalar product to raise the index, then
another one to lower it again to compute the scalar product between the two: it is ``a violation of the scalar product''
because one uses the structure and then its inverse. We do not need any: the statement is independent of the scalar product
\ts{5}{1}{1:39:16}.
Exercises on the sheet; all will be clear once we have a scalar product on every tangent space \ts{5}{1}{1:42:23}.
\end{aside}

\begin{example}[Counterexample: second derivatives are not a $(0,2)$ tensor field]\label{ex:tf-hessian}
Try to define components of a would-be $(0,2)$ tensor field by
\[
  H_{(x)\,ab}:=\partial_a\partial_b\bigl(f\circ x^{-1}\bigr)\qquad(\text{evaluated at }x(p)):
\]
two lower indices, a good candidate \ts{5}{1}{1:43:59}. Compute in the chart $y$ with
$f\circ y^{-1}=(f\circ x^{-1})\circ(x\circ y^{-1})$: differentiate first, evaluate at the end; the chain rule for the first
derivative and then the product rule and chain rule for the second give
\begin{align*}
  H_{(y)\,ab}\bigl(y(p)\bigr)
  &=\frac{\partial x^m}{\partial y^a}\frac{\partial x^n}{\partial y^b}\,H_{(x)\,mn}\bigl(x(p)\bigr)\\
  &\quad+\;\frac{\partial f}{\partial x^m}\Big|_{x(p)}\;\frac{\partial^2x^m}{\partial y^a\partial y^b}\Big|_{y(p)},
\end{align*}
\ts{5}{1}{1:44:42}.
The first term is the correct $(0,2)$ law; if the second term were absent it would be a tensor field, but it is present and
in general does not vanish: \emph{an extra term}. So the matrix of second derivatives of $f$ (the Jacobian-type matrix) is
\textbf{not} a $(0,2)$ tensor field: ``so fast can it go'' \ts{5}{1}{1:51:30}.
\end{example}
\begin{coursenote}[Second term: the Hessian transforms with a correction]
Since $\partial f/\partial x^m=(\mathrm{grad}f)_m$ and $\partial^2x^m/\partial y^a\partial y^b$ is the second derivative of the transition
map, the extra term is built from the gradient and from the transition map, not from $H$: it vanishes for linear changes of chart
(e.g.\ between Cartesian systems), but not for curvilinear ones. (Added; not said in class.)
\end{coursenote}
\paragraph{Outlook.} Next time we introduce a new structure that allows us to complete the $\ddot q$ of \cref{ex:tf-second}
to a vector, and in the same way the second derivative above to a $(0,2)$ tensor field: ``more structure'' (the linear connection)
\ts{5}{1}{1:52:54}.

% ==================================================================
\section*{Key concepts}
\begin{keyconcepts}
\textbf{Retrieval questions.}
\begin{enumerate}
  \item State the transformation law for vector components under a chart change; where does it come from?
  \item Why is $X^i\partial_i$ with arbitrary components not yet a vector? What is the demand?
  \item Show that $\dot q^i$ transforms as a vector and $\ddot q^i$ does not. What is the extra term?
  \item Define a vector field and a smooth vector field. What changes in the transformation law compared with a vector?
  \item Define an $(r,s)$ tensor field in a chart; how many components; what is the law for upper and lower indices?
  \item Why are the two Jacobians mutually inverse? What is $\delta^a{}_b$?
  \item Why is $\partial_jf$ a covector field and not a vector field? Why is $\partial_a\partial_bf$ not a $(0,2)$ tensor field?
  \item Why does ``the gradient is perpendicular to the level lines'' not need a scalar product when stated with covectors?
\end{enumerate}
\textbf{Concepts.}
\begin{itemize}
  \item Vector law (\cref{def:tf-vector}); constructing components and checking their law (\cref{ex:tf-first,ex:tf-second}).
  \item Vector fields (\cref{def:tf-vectorfield}); tensor fields and their law (\cref{def:tf-tensorfield,thm:tf-tensor-trafo}).
  \item Inverse Jacobians, Kronecker delta (\cref{thm:tf-inverse}); gradient as $(0,1)$ field, second derivatives not tensorial
        (\cref{ex:tf-gradient,ex:tf-hessian}).
\end{itemize}
\end{keyconcepts}

\flushside
\end{paracol}
